\section{DSL Requirements}

\subsection{General Language Rules}

Research indicates that children consistently struggle with strict syntax requirements, punctuation rules, language formalities, verbose statements, and inadequate debugging feedback \cite{klimekova_is_2015}.
Python addresses these challenges through its minimal syntax and readable structure, which explains its widespread adoption in introductory programming courses across universities worldwide \cite{klimekova_is_2015}.
These findings informed the decision to adopt Python as the primary source of inspiration for the DSL's syntax and overall design.

Empirical research on identifier naming conventions demonstrates that different styles --- camelCase, PascalCase, and snake\_case --- impose measurably different levels of visual effort and affect the speed at which identifiers are recognised by the reader. 
Studies using eye-tracking methodology found that snake\_case yields the lowest visual effort and the highest recognition speed, particularly among novice programmers \cite{sharif_eye_2010}. 
Consequently, snake\_case was selected as the identifier style for all predefined variables and functions in the DSL.

\subsection{Language Paradigm}

Given that the DSL targets children with minimal programming experience, an imperative paradigm was selected as the foundational approach. 
Empirical research demonstrates that imperative programming exhibits one of the lowest learning curves among major paradigms \cite{ugwumba_programming_2026}. 
This paradigm is particularly suitable for novice learners, as it aligns with sequential, step-by-step reasoning patterns and facilitates code comprehension during debugging processes.

While object-oriented programming presents comparable learning difficulty to imperative programming \cite{ugwumba_programming_2026}, it was deemed less appropriate for this context. 
The object-oriented paradigm emphasizes abstraction and code encapsulation, concepts that may prove challenging for the target demographic. 
Additionally, object-oriented approaches typically require more verbose code structures, introducing unnecessary complexity beyond the DSL's intended scope.

\subsection{Core Language Constructs}

The DSL should include variables to familiarize children with a core component of most programming languages. 
Variables are essential for storing coordinates, object IDs, and other Minecraft elements that users need to manipulate.

To enable process automation, the DSL will incorporate conditional statements and iterative constructs for handling repetitive tasks. 
As these represent fundamental components of most programming languages, their inclusion supports progressive learning objectives.

\subsection{Abstraction Level and Complexity of Functionalities}

To preserve user flexibility and support progressive skill development, the DSL will adopt a layered abstraction approach, providing both primitive and composite operations. 
Primitive operations include atomic actions such as block placement, directional orientation, jumping, clicking, and inventory management. 
Composite operations encompass higher-level tasks including wall construction, water removal from defined areas, and automated crop cultivation.

This architectural design enables users to construct custom high-level functionalities by composing primitive operations when predefined composite operations do not satisfy specific requirements. 
Such an approach promotes understanding of functional composition in computational thinking.

\subsection{DSL Limitations}

To prevent implementation complexity and breaking Minecraft official policies, the DSL will not change the game's state with means external to the user's in-game character. 
Thus, the DSL allows the user to interact with the game only through their character.

Although the usage of the DSL is intended for children, the installation process of the DSL requires external assistance.